
\documentclass{article}
\usepackage{colortbl}
\usepackage{makecell}
\usepackage{multirow}
\usepackage{supertabular}

\begin{document}

\newcounter{utterance}

\centering \large Interaction Transcript for game `eqbench', experiment `eq\_bench', episode 13 with reflexion{-}Qwen3.5{-}9B.
\vspace{24pt}

{ \footnotesize  \setcounter{utterance}{1}
\setlength{\tabcolsep}{0pt}
\begin{supertabular}{c@{$\;$}|p{.15\linewidth}@{}p{.15\linewidth}p{.15\linewidth}p{.15\linewidth}p{.15\linewidth}p{.15\linewidth}}
    \# & \multicolumn{2}{c}{Player} && \multicolumn{2}{c}{Game Master} \\
    \hline

    \theutterance \stepcounter{utterance}  
    & & & \multicolumn{4}{p{0.6\linewidth}}{
        \cellcolor[rgb]{0.9,0.9,0.9}{
            \makecell[{{p{\linewidth}}}]{
                \texttt{\tiny{[P1$\langle$GM]}}
                \texttt{You have attempted this same task before and did not succeed. These are the plans you wrote after those attempts.} \\
\texttt{Trial \#1: New plan: I failed because I assigned a high score to "Exasperation" (9) and low scores to others, but the dialogue shows Lorelei is not frustrated; she is consistently calm, logical, and reframing Harold's points rather than reacting with anger. Her tone is firm but not heated, making "Exasperation" an incorrect primary emotion. Next time, I will analyze the emotional valence of her responses: since she is trying to educate and persuade without raising her voice, "Amusement" is unlikely, but "Acceptance" of the conflict's reality is also low because she is actively fighting back. The most accurate emotion is "Dismay" or a similar sense of concern/frustration at the *situation* (the election's stakes) rather than Harold personally, but given the options, "Dismay" fits her worried tone about the future better than "Exasperation." I will re{-}evaluate by noting her use of phrases like "That's not the point" and "we can't afford not to act" as indicators of deep concern (Dismay) rather than annoyance, and ensure at least one score is significantly higher than the others to reflect the dominant mood. Specifically, I will lower Exasperation to 2{-}3 and raise Dismay to 8{-}9.} \\
\\ 
\\ 
\\ 
\texttt{Your task is to predict the likely emotional responses of a character in this dialogue:} \\
\\ 
\texttt{Harold: I don't understand why you're so worked up about this election. It's not like the world is going to end if your candidate doesn't win.} \\
\texttt{Lorelei: That's not the point, Harold. It's about what this election represents. It's about the future of our planet, our society.} \\
\texttt{Harold: And you think your candidate has all the answers? } \\
\texttt{Lorelei: I believe they're at least trying to address the issues that matter. Climate change, social justice, healthcare...} \\
\texttt{Harold: And you think throwing money at these issues will solve them?} \\
\texttt{Lorelei: It's not about throwing money, Harold. It's about investing in solutions, in our future.} \\
\texttt{Harold: And who's going to pay for all these grand plans? Us, the taxpayers. } \\
\texttt{Lorelei: Isn't that the point of taxes? To fund public services and infrastructure?} \\
\texttt{Harold: Yes, but there's a limit. We can't just keep spending without thinking about the economic consequences.} \\
\texttt{Lorelei: And we can't keep ignoring the environmental and social consequences of our actions either.} \\
\texttt{Harold: I'm not suggesting we ignore them, Lorelei. I'm just saying we need to be realistic about what we can afford.} \\
\texttt{Lorelei: And I'm saying we can't afford not to act. The cost of inaction is far greater.} \\
\texttt{Harold: And who decides what action is necessary? The government? The same government that can't even balance a budget?} \\
\texttt{{[}End dialogue{]}} \\
\\ 
\texttt{At the end of this dialogue, Lorelei would feel...} \\
\texttt{Dismay} \\
\texttt{Exasperation} \\
\texttt{Amusement} \\
\texttt{Acceptance} \\
\\ 
\texttt{Give each of these possible emotions a score from 0{-}10 for the relative intensity that they are likely to be feeling each.} \\
\\ 
\texttt{You must output in the following format, including headings (of course, you should give your own scores), with no additional commentary:} \\
\\ 
\texttt{Dismay: <score>} \\
\texttt{Exasperation: <score>} \\
\texttt{Amusement: <score>} \\
\texttt{Acceptance: <score>} \\
\\ 
\\ 
\texttt{{[}End of answer{]}} \\
\\ 
\texttt{Remember: zero is a valid score, meaning they are likely not feeling that emotion. You must score at least one emotion > 0.} \\
\\ 
\texttt{Your answer:} \\
            }
        }
    }
     \\ \\

    \theutterance \stepcounter{utterance}  
    & \multicolumn{4}{p{0.6\linewidth}}{
        \cellcolor[rgb]{0.561,0.737,0.561}{
            \makecell[{{p{\linewidth}}}]{
                \texttt{\tiny{[P1$\rangle$GM]}}
                \texttt{Dismay: 9} \\
\texttt{Exasperation: 2} \\
\texttt{Amusement: 0} \\
\texttt{Acceptance: 1} \\
            }
        }
    }
    & & \\ \\

    \theutterance \stepcounter{utterance}  
    & & & \multicolumn{2}{p{0.3\linewidth}}{
        \cellcolor[rgb]{0.9,0.9,0.9}{
            \makecell[{{p{\linewidth}}}]{
                \texttt{\tiny{[GM$|$GM]}}
                \texttt{Dismay: 9} \\
\texttt{Exasperation: 2} \\
\texttt{Amusement: 0} \\
\texttt{Acceptance: 1} \\
            }
        }
    }
    & & \\ \\

    \theutterance \stepcounter{utterance}  
    & & & \multicolumn{2}{p{0.3\linewidth}}{
        \cellcolor[rgb]{0.9,0.9,0.9}{
            \makecell[{{p{\linewidth}}}]{
                \texttt{\tiny{[GM$|$GM]}}
                \texttt{Dismay: 4} \\
\texttt{Exasperation: 8} \\
\texttt{Amusement: 0} \\
\texttt{Acceptance: 0} \\
            }
        }
    }
    & & \\ \\

    \theutterance \stepcounter{utterance}  
    & & & \multicolumn{2}{p{0.3\linewidth}}{
        \cellcolor[rgb]{0.9,0.9,0.9}{
            \makecell[{{p{\linewidth}}}]{
                \texttt{\tiny{[GM$|$GM]}}
                \texttt{game\_result = LOSE} \\
            }
        }
    }
    & & \\ \\

    \theutterance \stepcounter{utterance}  
    & & \multicolumn{4}{p{0.6\linewidth}}{
        \cellcolor[rgb]{0.95,0.95,0.95}{
            \makecell[{{p{\linewidth}}}]{
                \texttt{\tiny{[P1$\leftrightarrow$P]}}
                \texttt{Post{-}episode Reflexion (not sent to GM):} \\
\texttt{New plan: I failed because I assigned a high score to "Exasperation" (9), misinterpreting Lorelei's firm, logical rebuttals as signs of frustration or annoyance. In reality, her tone is consistently calm, reasoned, and focused on the broader stakes (climate, society) rather than reacting emotionally to Harold's personal attacks. She is not "worked up"; she is steadfastly advocating for her principles. Therefore, "Exasperation" should be scored very low (0{-}2). The dominant emotion is actually a deep, worried concern about the future and the consequences of inaction, which aligns perfectly with "Dismay." "Amusement" is incorrect as the tone is serious, and "Acceptance" is incorrect because she is actively challenging Harold's skepticism, not passively accepting his view. My new strategy is to identify "Dismay" as the primary emotion (score 8{-}10) reflecting her anxiety about the planet's future, while keeping "Exasperation" near zero to reflect her lack of personal irritation, and ensuring the remaining scores reflect the absence of humor or resignation.} \\
            }
        }
    }
    & \\ \\

\end{supertabular}
}

\end{document}
